\subsubsection{Resolution Marker}
\label{SEC:ResolutionMarker}
\key*{}[marker!resolution marker]

%%%%%%%%%%%%%%%%%%%%%%%%%%%%%%%%%%%%%%%%%%%%%%%%%%%%%%%%%%%%%%%%%%%%%%%%%%%%%%%%

%! The idea behind the resolution marker is that numerical instabilities
%! often appear in the highest polynomial mode(s).

As discussed in \sref{SEC:Filtering}, numerical instabilities often
appear in the highest polynomial mode(s).
%
An easy way to distinguish between well-resolved and possibly unstable
cells is to examine the highest mode of the solution within each cell.
 
%%%%%%%%%%%%%%%%%%%%%%%%%%%%%%%%%%%%%%%%%%%%%%%%%%%%%%%%%%%%%%%%%%%%%%%%%%%%%%%%

Recall that the \key{modal form}[polynomial!modal form] of the
solution polynomial for a given flow variable is given by \eref{EQN:Modal_Form}
\begin{align*}
       u (r) =& \sum_{n=1}^{N_p} \uModal[n] (t) \psi_n (r)
& \psi_n (r) =& r^{n-1}
\end{align*}
where $\uModal[n]$ represents the modal coefficients, with $\uModal[1]$ being
the coefficient for the lowest\slash constant mode and $\uModal[N_p]$ being
the coefficient for the highest\slash degree-\PO{} mode.
%
A diagonal filter matrix can be applied to remove the highest mode of the
solution.
\begin{align}
    \bm{\underline{\uModal[]}}
 =& \bm{\sigma} \bm{\uModal[]}
  &
    \sigma_n
 =& \begin{cases}
      0 & n = N_p \\
      1 & \text{otherwise}
    \end{cases}
\end{align}
 
A comparison of this filtered solution to the actual solution will
provide a test for the possibility of numerical instabilities in
the cell.
%
Namely, in a well-resolved solution,
$\bm{\uModal[]} \simeq \bm{\underline{\uModal[]}}$.
%
Any cells in which this comparison shows a difference that is beyond
a certain threshold are marked as ``troubled'' cells.

%%%%%%%%%%%%%%%%%%%%%%%%%%%%%%%%%%%%%%%%%%%%%%%%%%%%%%%%%%%%%%%%%%%%%%%%%%%%%%%%

To do this, we compute the local resolution marker
as\cite{Persson_2006,Barter_2007,Barter_2008}
\begin{align}
    \label{EQN:ResolutionMarker}
    \marker \left( \bm{\uModal[]} , \bm{\underline{\uModal[]}} \right)
 =& \log_{10}
    \left[
      \frac{
        \left( \bm{\uModal[]} - \bm{\underline{\uModal[]}} \right) \cdot
        \left( \bm{\uModal[]} - \bm{\underline{\uModal[]}} \right)
      }{
               \bm{\uModal[]} \cdot        \bm{\uModal[]}
      }
    \right]
\end{align}
where $\bm{\uModal[]}$ is the user-specified solution variable
to check (density is recommended).
%
If all of the conserved variables are specified, then
\eref{EQN:ResolutionMarker} is computed for each one and the
maximum value is used.
%
Note that the numerator represents the norm of the highest mode
and the denominator represents the norm of all the modes combined.
%
The marker will have values in the range $\marker = [-\infty,0]$ where
the lower value indicates that nothing is contained in the highest mode and 
the upper value indicates that everything is contained in the highest mode.

%%%%%%%%%%%%%%%%%%%%%%%%%%%%%%%%%%%%%%%%%%%%%%%%%%%%%%%%%%%%%%%%%%%%%%%%%%%%%%%%

Next, we compute the local scaled marker as described in \sref{SEC:ScaledMarker}.

%%%%%%%%%%%%%%%%%%%%%%%%%%%%%%%%%%%%%%%%%%%%%%%%%%%%%%%%%%%%%%%%%%%%%%%%%%%%%%%%

%!% AIAA-2022-1813 by Spiegel used:
%!\begin{align}
%!    {\cal S} >& - \left( \psi + 4.25 \log_{10} \left( \PO \right) \right) &
%!\end{align}
%!where $\psi$ is a user-specified constant that will depend on the
%!problem and the polynomial order $\PO$.

%%%%%%%%%%%%%%%%%%%%%%%%%%%%%%%%%%%%%%%%%%%%%%%%%%%%%%%%%%%%%%%%%%%%%%%%%%%%%%%%

\SeeAlso{
\inp[nmlLimiter::MarkerOption]{Marker\_Option},
\inp[nmlLimiter::MarkerVariable]{Marker\_Variable}
}


%%%%%%%%%%%%%%%%%%%%%%%%%%%%%%%%%%%%%%%%%%%%%%%%%%%%%%%%%%%%%%%%%%%%%%%%%%%%%%%%

